%
% 2-kalkuel.tex
%
% (c) 2026 Prof Dr Andreas Müller
%
\section{Der Differentialkalkül}
\kopfrechts{Der Differentialkalkül}
Zur Zeit der Entstehung der Infinitesimalrechnung war die numerische
Berechnung meistens eine zu aufwendige Option.
An erster Stelle stand wann immer möglich die algebraische Rechnung.
Erst nach allen möglichen algebraischen Vereinfachungen war daran zu
denkten, konkrete Zahlenwerte einzusetzen und numerische Resultate
zu vergleichen.
Die Infinitesimalrechnung entstand daher zuerst primär als ein Kalkül,
was sich auch in dem im englischsprachigen Raum üblichen Namen
\emph{Calculus} widerspiegelt.

%
% Eulers Einführung in die Analysis
%
\subsection{Eulers Einführung in die Analysis}
Das Buch \emph{Introductio in analysin infinitorum}
\cite{buch:introductio1,buch:introductio2}
beginnt Leonhard Euler, den Leser in die Ideen der Inifinitesimalrechnung
einzuführen.
\index{Euler, Leonhard}%
Es behandelt alle Aspekte der Infinitesimalrechnung, die ohne das
Konzept der Ableitung eingführt werden können.
Es behandelt im ersten Band Funktionen einer und mehrere Variablen,
die trigonometrischen Funktionen, Reihen, Produktdarstellungen,
\index{trigonometrische Funktion}%
\index{Reihe}%
\index{Produktdarstellung}%
Kettenbrüche und Anwendungen dieser Themen.
\index{Kettenbruch}%
Man findet darin auch die Formel $e^{i\varphi}=\cos\varphi+i\sin\varphi$.
Im zweiten Band \cite{buch:introductio2} behandelt er sowohl 
algebraisch definierte Kurven wie auch Graphen transzendenter
\index{Kurve}%
\index{Graph}%
Funktionen.

Erst in seinen späteren Werk \emph{Institutiones calculi differentialis}
\cite{buch:institutiones} behandelt er die Ableitung und das Integral.
Anwendungen die wie Bestimmung von Extrema oder Differentialgleichungen
\index{Extrema}%
werden darin ausführlich mit vielen Besipielen durchgerechnet.
Das Integralzeichen, wie es Leibniz eingeführt hat, kommt jedoch im
\index{Integralzeichen}%
ganzen Buch nicht vor.
Bei der Lösung von Differentialgleichungen bedient er sich der algebraischen
\index{Differentialgleichung}%
Regeln für die Ableitung, ohne je einen formellen Integralkalkül
entwickelt zu haben.

Eulers Werk ist typische für die im 18.~Jahrhundert vorherrschende
Betrachtungsweise der Infinitesimalrechnung.
Das Problem der Definition einer Funktion und der Begriffe Grenzwert
und Ableitung werden eher intuitiv behandelt.
Mit den ``unendlich kleinen'' Grössen $dx$ und $dy$ wird nach den
üblichen Regeln der Algebra gerechnet, eine Konstruktion eines Integrals
wird nicht durchgeführt. 
Sogar Differentialgleichungen sind nicht Gleichungen zwischen Ableitungen
einer gesuchten Funktion, sondern tatsächlich Gleichungen zwischen
den Differentialen $dx$, $dy$, $d^2y$ und $dx^2$.

%
% Die Rechenregeln der klassischen Infinitesimalrechnung
%
\subsection{Die Rechenregeln der klassischen Infinitesimalrechnung}
Für die praktische Arbeit mit Ableitungen im Stile Eulers genügt es,
die folgenden Regeln zu kennen.
Regeln für Funktionen:
\begin{align}
&\text{Linearität:}&
\frac{d}{dx}(f(x)\pm g(x))
&=
\frac{d}{dx}f(x)
\pm
\frac{d}{dx}g(x)
\label{buch:einleitung:kalkuel:eqn:lin1}
\\
&&
\frac{d}{dx}(\alpha f(x))
&=
\alpha \frac{d}{dx}f(x)
\label{buch:einleitung:kalkuel:eqn:lin2}
\\
&\text{Produkt:}&
\frac{d}{dx}\bigl(f(x)g(x)\bigr)
&=
\biggl(\frac{d}{dx}f(x)\biggr)
g(x)
+
f(x)
\biggl(\frac{d}{dx}g(x)\biggr).
\label{buch:einleitung:kalkuel:eqn:prod}
\intertext{\index{linear}\index{Produktregel}%
Ein Regel, die Konstanten charakterisiert:}
&\text{Konstanten:}&
\frac{d}{dx}c&=0\qquad\text{wenn $c$ eine Konstante ist.}
\label{buch:einleitung:kalkuel:eqn:const}
\intertext{Eine Regel, die die unabhängige Variable auszeichnet:}
&\text{Variable:}&
\frac{d}{dx}x&=1.
\label{buch:einleitung:kalkuel:eqn:var}
\intertext{Die Ableitung des Reziproken $1/f(x)$ lässt sich daraus und aus
der Eigenschaft $f(x)\cdot 1/f(x)=1$ mit der Produktregel gewinnen,
\index{Quotientenregel}%
ebenso die Quotientenregel}
&\text{Quotienten:}&
\frac{d}{dx} \frac{f(x)}{g(x)}
&=
\frac{\frac{df}{dx}(x)g(x) - f(x)\frac{dg}{dx}(x)}{g(x)^2}.
\intertext{Dazu kommt noch die Kettenregel}
&\text{Kettenregel:}&
\frac{d}{dx} f(g(x))&=\frac{df}{dx}(g(x))\frac{dg}{dx}(x),
\end{align}
\index{Kettenregel}%
die für zusammengesetzte Funktionen nützlich ist.
\index{Zusammensetzung}

In der klassischen Entwicklung der Infinitesimalrechnung werden diese
Regeln aus der Definition der Ableitung als Grenzewert von
Differenzenquotienten gewonnen.
\index{Differenzenquotient}%
Man kann sich aber auch auf den Standpunkt stellen, dass diese Regeln
einen speziellen Operator auf allen Funktionen einer geeignet definierten
Funktionenmenge charakterisieren.
Die Funktionenmenge muss alle arithmetischen Operationen gestatten,
man muss also Funktionen addieren, subtrahieren, multiplizieren und
dividieren können.
Funktionen können mit Konstanten multipliziert werden.
Es muss möglich sein, Gleichungen mit Funktionsausdrücken aufzulösen.
Diese Eigenschaften lassen sich zusammenfassen in die Aussage, dass
die Menge $F$ der Funktionen folgende Eigenschaften hat:
\begin{itemize}
\item
$F$ ist ein Körper, d.~h.~alle vier Grundoperationen sind immer ausführbar
ausser die Division durch die Konstanten $0$.
\item
Sie enthält einen Teilkörper $K\subset F$ von Konstanten, der mindestens
die rationalen Zahlen enthält $\mathbb{Q}\subset F$.
$K$ heisst auch der Konstantenkörper.
Meistens geht man von den reellen Zahlen $\mathbb{R}$ als Zahlenkörper aus.
\index{Konstantenkorper@Konstantenkörper}%
Für die Computeralgebra wird jedoch der Körper $\mathbb{Q}$ vorgezogen.
\item
$F$ ist eine Algebra über dem Konstantenkörper $K$.
Da $K\subset F$ ist, ist dies automatisch erfüllt.
Es ist aber nützlich, diese Eigenschaft explizit zu erwähnen, weil die
Definition einer Derivation nur die Algebrastruktur braucht.
\index{Derivation}%
\end{itemize}
Man beachte, dass diese Struktur das Konzept der Ableitung noch nicht
definiert.
Der Funktionenkörper ist im Wesentlichen der Schauplatz von Eulers
\index{Funktionenkorper@Funktionenkörper}%
\emph{Introduction}
\cite{buch:introductio1,buch:introductio2}.

Die Ableitung ist ein zusätzlicher Operator $D$, der auf Funktionen in
$F$ wirkt.
Die beiden Regeln \eqref{buch:einleitung:kalkuel:eqn:lin1} und
\eqref{buch:einleitung:kalkuel:eqn:lin2} bekommen die Form
\begin{align*}
 D(f+g) &= Df + Dg\\
 D(\alpha f) &= \alpha Df
\intertext{für $f,g\in F$ und $\alpha\in K$, die besagt,
dass $D$ ein linearer Operator ist.
\index{linearer Operator}%
Aus der Produktregel \eqref{buch:einleitung:kalkuel:eqn:prod}
\index{Produktregel}%
wird}
D(fg) &= (Df)g+f(Dg)
\intertext{für $f,g\in F$.
Diese Eigenschaft von $D$ heisst auch \emph{Derivation}.
\index{Derivation}%
Schliesslich identifiziert
\eqref{buch:einleitung:kalkuel:eqn:const}
die Konstanten und
\eqref{buch:einleitung:kalkuel:eqn:var}
wird zu}
Dx&=1,
\end{align*}
was die unabhängige Variable $x$ charakterisiert, nach der der Operator
$D$ ableitet.
Die Vorgabe eines linearen Operators auf einer Algebra über den
rationalen Zahlen mit obigen Eigenschaften ist alles, was man braucht, um 
die Rechnungen der Infinitesimalrechnung durchführen zu können.
Sie entpuppt sich somit als hauptsächlich algebraischer Kalkül.

Im algebraischen Kalkül spielt die Kettenregel kein Rolle, da in einem
\index{Kettenregel}%
Funktionenkörper das Konzept der Zusammensetzung von Funktionen gar
nicht definiert ist.
Dies beschränkt die Möglichkeiten, Ableitungen in diesem Körper zu berechnen,
jedoch nicht.

